% =====================================================================
\chapter{Related work}\label{ch:related}
% =====================================================================
\paragraph{The base paper and its citers.}
Tan, Zheng and Vijayakumar~\cite{tan2023} are, to their knowledge, the first to use a tethered UAV for authentication in infrastructure-less IoVs. We
screened the 94 works that cite it; none reports an attack on its CRT key distribution or on its key-escrow claim, and none protects the tether itself.

\paragraph{CRT and polynomial group-key distribution.}
CRT-based broadcasting goes back to the ``secure lock'' of Chiou and Chen~\cite{chiou89}. Sharma and Purushothama~\cite{sp24} show that prominent
CRT-based group-key schemes fail to keep backward secrecy when a user is corrupted; our insider results (F2a, F2b) are an explicit instance of this
weakness for the base scheme. Masked-polynomial constructions are known to be fragile in other settings too: Albrecht, Gentry, Halevi and
Katz~\cite{agh09} break key-predistribution schemes based on ``perturbation polynomials''. Our outsider attack (F2c) differs from these: it needs no
corrupted user, only two public broadcasts, and it rests on Vieta's formulas. Logical key hierarchies~\cite{wgl98} are the standard way to rekey large
groups with per-member delivery, and C1 can use them.

\paragraph{Certificateless signatures and batch verification.}
Certificateless public-key cryptography was introduced by Al-Riyami and Paterson~\cite{alriyami} to remove both certificates and key escrow; Girault's
trust levels~\cite{girault} grade how much an authority can still cheat. Our credential is the Schnorr-based identity construction of Galindo and
Garcia~\cite{gg09}, whose security proof was revisited by Chatterjee, Kamath and Kumar~\cite{ckk12}, extended with the user's own secret. Schnorr
signatures~\cite{schnorr} and the forking lemma~\cite{forking,bnforking} give the security arguments, and the small-exponent test of Bellare, Garay and
Rabin~\cite{bgr} gives sound batch verification.

\paragraph{Handover and tickets.}
Ticket-based authentication is classic (Kerberos~\cite{kerberos94}). Handover schemes for UAV and vehicular networks typically assume an authority that
remains reachable, or rely on physical unclonable functions or a blockchain. C2 works with the authority offline and binds the ticket to signatures in both
directions, which matters precisely when the issuing TUAV may be captured.

\paragraph{Byzantine fault tolerance.}
C3 borrows the quorum-certificate pattern of PBFT's view change~\cite{castro99}; because eviction is monotone, a quorum of $f+1$ suffices for safety.

\paragraph{Mediated and threshold signatures.}
Mediated cryptography~\cite{bdtw01} splits a key between a user and a security mediator that can revoke it instantly; Chow, Boyd and Gonz\'alez
Nieto~\cite{cbn06} combine it with certificateless keys, and Yap et al.~\cite{yap07} give pairing-free mediated certificateless signatures. Two-party and
threshold Schnorr signing is standardised as FROST~\cite{frost20,rfc9591}, and the ROS attack~\cite{ros21} shows why naive concurrent signing is unsafe.
C4 applies these ideas to a new place: the mediator is the ECRV at the other end of the tether, and the revocation trigger is physical.

\paragraph{Formal verification.}
We use the symbolic Dolev--Yao model~\cite{dolevyao} and the ProVerif verifier~\cite{blanchet16} for the handover protocol.
